% Setting up the document and importing necessary packages
\documentclass[12pt,a4paper]{article}

% Chinese support packages
\usepackage{xeCJK}
\usepackage{fontspec}
\setCJKmainfont{Noto Serif CJK TC}
\setCJKsansfont{Noto Sans CJK TC}
\setCJKmonofont{Noto Sans CJK TC}

% Other necessary packages
\usepackage{geometry}
\usepackage{titlesec}
\usepackage{enumitem}
\usepackage{color}
\usepackage{colortbl}
\usepackage{tcolorbox}
\usepackage{booktabs}
\usepackage{longtable}
\usepackage{array}
\usepackage{graphicx}
\usepackage{tikz}
\usepackage{mdframed}
\usepackage{fancyhdr}
\usepackage{hyperref}
\usepackage{multicol}
\usepackage{datetime2}
\usepackage{natbib}

\hypersetup{
    colorlinks=true,
    linkcolor=onco_teal,
    citecolor=onco_teal,
    urlcolor=onco_teal,
    pdfborder={0 0 0}
}

% Page setup
\geometry{
    top=2.5cm,
    bottom=2.5cm,
    left=2cm,
    right=2cm,
    headheight=15pt
}

% Color definitions (Oncology themed colors)
\definecolor{onco_teal}{RGB}{0,128,128}
\definecolor{onco_gold}{RGB}{218,165,32}
\definecolor{onco_gray}{RGB}{120,120,120}
\definecolor{highlight_yellow}{RGB}{255,250,205}

% Title format settings
\titleformat{\section}[block]{\Large\bfseries\color{onco_teal}}{\thesection}{10pt}{}
\titleformat{\subsection}[block]{\large\bfseries\color{onco_gold}}{\thesubsection}{8pt}{}
\titleformat{\subsubsection}[block]{\normalsize\bfseries\color{onco_gray}}{\thesubsubsection}{6pt}{}

% Custom environment
\newmdenv[
    linecolor=onco_teal,
    backgroundcolor=highlight_yellow,
    linewidth=2pt,
    topline=true,
    bottomline=true,
    leftline=true,
    rightline=true,
    innertopmargin=10pt,
    innerbottommargin=10pt,
    innerleftmargin=10pt,
    innerrightmargin=10pt
]{importantbox}

\newcommand{\note}[1]{
    \par
    \noindent
    \begin{tcolorbox}[
        colback=onco_gold!5!white,
        colframe=onco_gold,
        title=\textbf{重點提醒},
        fonttitle=\bfseries,
        boxrule=0.5pt,
        arc=3pt,
        left=6pt,
        right=6pt,
        top=6pt,
        bottom=6pt
    ]
    #1
    \end{tcolorbox}
}

% Header and footer settings
\pagestyle{fancy}
\fancyhf{}
\fancyhead[L]{\textcolor{onco_teal}{內科部住院醫師教學文件}}
\fancyhead[R]{\textcolor{onco_gold}{腫瘤科EPAs}}
\fancyfoot[C]{\textcolor{onco_gray}{\thepage}}
\renewcommand{\headrulewidth}{1pt}
\renewcommand{\headrule}{\hbox to\headwidth{\color{onco_teal}\leaders\hrule height \headrulewidth\hfill}}

% Date format
\DTMsetstyle{yyyy-mm-dd}
\newcommand{\mydate}{\DTMtoday}

% Document start
\begin{document}

% Title page
\begin{titlepage}
    \centering
    {\LARGE\textcolor{onco_teal}{\textbf{內科部住院醫師教學文件}}\par}
    \vspace{1cm}
    {\huge\bfseries 腫瘤科可信賴專業活動EPAs\par}
    \vspace{0.5cm}
    {\Large\textcolor{onco_gold}{\textit{Oncology Entrustable Professional Activities}}\par}
    \vspace{1cm}
    
    \begin{tcolorbox}[
        colback=onco_gold!5!white,
        colframe=onco_gold,
        width=0.8\textwidth,
        arc=10pt,
        boxrule=1pt
    ]
    \centering
    {\large\textbf{目標對象}\par}
    \vspace{0.3cm}
    內科部住院醫師R1-R3\par
    腫瘤科專科訓練醫師\par
    \vspace{0.5cm}
    {\large\textbf{文件版本}\par}
    \vspace{0.3cm}
    第1.0版(10個EPAs)\par
    發布日期:\mydate\par
    \end{tcolorbox}

    \vfill
    
    {\large 腫瘤科 \\
    適用於CCC評量系統\par}
    
    \vspace{0.5cm}
\end{titlepage}

\newpage
\tableofcontents
\newpage

% Main content
\section{腫瘤科EPAs概述}

\subsection{EPA定義}
Entrustable Professional Activities (EPA) 是指可信賴委託給住院醫師執行的專業活動單元。腫瘤科EPAs涵蓋了癌症診斷、治療與照護的核心能力,是所有腫瘤科醫師必須精熟的專業技能。

\vspace{0.5cm}
\note{
腫瘤科EPAs評量著重於能否安全、穩定地執行完整的癌症診療流程,並能整合資訊做出適當臨床決策。每個EPA都包含明確的里程碑發展階段,從初學者到專家的完整能力建構。本版本包含10個核心EPA,從基礎的腫瘤學病史與病理判讀,到複雜的化療管理、標靶治療與緩和照護,提供完整的腫瘤疾病診療能力架構。
}

\subsection{學習目標}
完成腫瘤科EPAs後,住院醫師應能夠:
\begin{itemize}[nosep]
    \item 熟練進行腫瘤學病史詢問與理學檢查。
    \item 判讀病理報告與分子標記。
    \item 執行癌症分期與預後評估。
    \item 選擇與執行化學治療方案。
    \item 使用標靶治療與免疫治療。
    \item 評估放射治療適應症與協調。
    \item 診斷與處理腫瘤急症。
    \item 管理化療副作用與支持性照護。
    \item 執行癌症疼痛控制與緩和醫療。
    \item 協調腫瘤病患跨科部整合照護。
\end{itemize}

\subsection{委託等級}
\begin{table}[h]
    \centering
    \caption{EPA委託等級定義}
    \begin{tabular}{p{2cm} p{3cm} p{9cm}}
        \toprule
        \textbf{等級} & \textbf{委託程度} & \textbf{描述} \\
        \midrule
        Level 1 & 觀察學習 & 僅能觀察他人執行,無法獨立操作 \\
        Level 2 & 直接監督 & 可在主治醫師直接監督下執行 \\
        Level 3 & 間接監督 & 可獨立執行,但需回報並接受指導 \\
        Level 4 & 監督可得 & 可獨立執行,監督者隨時可得 \\
        Level 5 & 完全獨立 & 可獨立執行並指導他人 \\
        \bottomrule
    \end{tabular}
\end{table}

\newpage
\section{腫瘤科10大EPAs詳述}

\begin{longtable}{|p{1.2cm}|p{3.8cm}|p{8.5cm}|}
\caption{腫瘤科EPAs完整架構} \\
\hline
\rowcolor{onco_gold!20}
\textbf{EPA} & \textbf{專業活動名稱} & \textbf{核心能力要素與里程碑} \\
\hline
\endfirsthead

\multicolumn{3}{c}%
{{\bfseries \tablename\ \thetable{} -- 接續上頁}} \\
\hline
\rowcolor{onco_gold!20}
\textbf{EPA} & \textbf{專業活動名稱} & \textbf{核心能力要素與里程碑} \\
\hline
\endhead

\hline \multicolumn{3}{|r|}{{接續下頁}} \\ \hline
\endfoot

\hline \hline
\endlastfoot

\rowcolor{onco_teal!10}
\multicolumn{3}{|c|}{\textbf{EPA 1: 腫瘤學病史詢問與理學檢查}} \\
\hline
\textbf{描述} & \multicolumn{2}{|p{12.5cm}|}{能夠獨立進行癌症相關的完整病史詢問與理學檢查,包括B症狀評估、腫瘤家族史、環境暴露史,並能初步評估疾病範圍與功能狀態。} \\
\hline
\textbf{核心能力} & 警示症狀 & 體重減輕、不明發燒、盜汗、腫塊、出血 \\
\cline{2-3}
& 器官特定症狀 & 咳血、吞嚥困難、腸阻塞、黃疸評估 \\
\cline{2-3}
& 家族史 & 遺傳性癌症症候群、家族癌症史 \\
\cline{2-3}
& 環境暴露 & 吸菸、職業暴露、放射線、致癌物接觸 \\
\cline{2-3}
& 理學檢查 & 淋巴結、腫塊、肝脾、神經學檢查 \\
\cline{2-3}
& 功能評估 & ECOG PS、Karnofsky score評估 \\
\hline
\textbf{Level 1} & Novice & 能詢問基本腫瘤症狀,需要監導者引導完成 \\
\hline
\textbf{Level 2} & Advanced Beginner & 能系統性完成腫瘤病史詢問,執行基本理學檢查 \\
\hline
\textbf{Level 3} & Competent & 能獨立完成複雜病史詢問,整合症狀進行初步診斷 \\
\hline
\textbf{Level 4} & Proficient & 能在困難情況下獲取準確病史,識別罕見癌症表現 \\
\hline
\textbf{Level 5} & Expert & 能指導他人進行病史詢問,建立系統性評估流程 \\
\hline

\rowcolor{onco_teal!10}
\multicolumn{3}{|c|}{\textbf{EPA 2: 病理報告與分子標記判讀}} \\
\hline
\textbf{描述} & \multicolumn{2}{|p{12.5cm}|}{能夠獨立判讀病理報告,理解組織型態、分化程度、侵犯範圍,並能判讀分子標記,整合資訊提出治療建議。} \\
\hline
\textbf{核心能力} & 病理診斷 & 組織型態、分化程度、特殊染色 \\
\cline{2-3}
& 腫瘤分級 & WHO分級、Gleason score、Nottingham grade \\
\cline{2-3}
& 侵犯評估 & 深度、血管侵犯、淋巴侵犯、切緣評估 \\
\cline{2-3}
& 免疫組化 & ER/PR/HER2、PD-L1、Ki-67判讀 \\
\cline{2-3}
& 分子檢測 & EGFR、ALK、KRAS、MSI、TMB判讀 \\
\cline{2-3}
& 治療意義 & 標記與治療選擇關聯性 \\
\hline
\textbf{Level 1} & Novice & 能閱讀基本病理報告,理解診斷名稱 \\
\hline
\textbf{Level 2} & Advanced Beginner & 能判讀常見病理資訊,理解分級意義 \\
\hline
\textbf{Level 3} & Competent & 能獨立判讀病理報告,整合分子標記 \\
\hline
\textbf{Level 4} & Proficient & 能判讀複雜病理,評估預後與治療意義 \\
\hline
\textbf{Level 5} & Expert & 成為病理判讀專家,能教導複雜判讀 \\
\hline

\rowcolor{onco_teal!10}
\multicolumn{3}{|c|}{\textbf{EPA 3: 癌症分期與預後評估}} \\
\hline
\textbf{描述} & \multicolumn{2}{|p{12.5cm}|}{能夠執行癌症分期評估,選擇適當檢查,判讀分期影像,應用TNM系統,並能評估預後因子與制定治療策略。} \\
\hline
\textbf{核心能力} & TNM分期 & AJCC分期系統應用、版本差異理解 \\
\cline{2-3}
& 影像評估 & CT、MRI、PET-CT在分期的應用 \\
\cline{2-3}
& 分期檢查 & 依癌別選擇適當分期檢查組合 \\
\cline{2-3}
& 預後評估 & 腫瘤大小、分化、轉移、分子標記整合 \\
\cline{2-3}
& 分期會議 & MDT討論、分期資訊呈現 \\
\cline{2-3}
& 治療決策 & 依分期選擇治療目標與策略 \\
\hline
\textbf{Level 1} & Novice & 能理解基本TNM概念,知道常見檢查 \\
\hline
\textbf{Level 2} & Advanced Beginner & 能執行標準分期評估,應用TNM分期 \\
\hline
\textbf{Level 3} & Competent & 能獨立完成分期評估,制定檢查計畫 \\
\hline
\textbf{Level 4} & Proficient & 能處理複雜分期,整合多重預後因子 \\
\hline
\textbf{Level 5} & Expert & 成為分期專家,領導MDT討論 \\
\hline

\rowcolor{onco_teal!10}
\multicolumn{3}{|c|}{\textbf{EPA 4: 化學治療方案選擇與執行}} \\
\hline
\textbf{描述} & \multicolumn{2}{|p{12.5cm}|}{能夠依據癌症類型、分期、病人狀況選擇適當化療方案,計算劑量,執行前評估,並監測治療反應。} \\
\hline
\textbf{核心能力} & 方案選擇 & 依癌別、分期選擇標準化療方案 \\
\cline{2-3}
& 劑量計算 & BSA計算、器官功能調整、老年劑量 \\
\cline{2-3}
& 治療前評估 & 血球、肝腎功能、心肺功能評估 \\
\cline{2-3}
& 給藥順序 & 藥物順序、時間間隔、水化需求 \\
\cline{2-3}
& 反應評估 & RECIST criteria、療效評估時機 \\
\cline{2-3}
& 療程調整 & 劑量減量、延遲、中止決策 \\
\hline
\textbf{Level 1} & Novice & 能了解化療基本概念,知道常見方案 \\
\hline
\textbf{Level 2} & Advanced Beginner & 能執行標準化療方案,計算基本劑量 \\
\hline
\textbf{Level 3} & Competent & 能獨立選擇與執行化療,監測反應 \\
\hline
\textbf{Level 4} & Proficient & 能處理複雜情況,調整化療策略 \\
\hline
\textbf{Level 5} & Expert & 成為化療專家,設計個人化方案 \\
\hline

\rowcolor{onco_teal!10}
\multicolumn{3}{|c|}{\textbf{EPA 5: 標靶治療與免疫治療}} \\
\hline
\textbf{描述} & \multicolumn{2}{|p{12.5cm}|}{能夠評估標靶治療與免疫治療適應症,選擇適當藥物,監測療效與特殊副作用,並能處理免疫相關不良反應。} \\
\hline
\textbf{核心能力} & 標靶藥物 & EGFR-TKI、ALK-i、HER2-i等適應症 \\
\cline{2-3}
& 生物標記 & 標靶治療生物標記檢測與判讀 \\
\cline{2-3}
& 免疫治療 & PD-1/PD-L1抑制劑、CTLA-4抑制劑 \\
\cline{2-3}
& irAE識別 & 免疫相關肺炎、腸炎、肝炎、內分泌 \\
\cline{2-3}
& irAE處理 & 類固醇使用、免疫抑制劑、停藥決策 \\
\cline{2-3}
& 抗藥性 & 抗藥機轉、後線治療選擇 \\
\hline
\textbf{Level 1} & Novice & 能了解標靶與免疫治療基本概念 \\
\hline
\textbf{Level 2} & Advanced Beginner & 能執行標準標靶治療,識別常見irAE \\
\hline
\textbf{Level 3} & Competent & 能獨立使用標靶與免疫治療 \\
\hline
\textbf{Level 4} & Proficient & 能處理複雜irAE,選擇後線治療 \\
\hline
\textbf{Level 5} & Expert & 成為標靶免疫專家,參與臨床試驗 \\
\hline

\rowcolor{onco_teal!10}
\multicolumn{3}{|c|}{\textbf{EPA 6: 放射治療適應症評估與協調}} \\
\hline
\textbf{描述} & \multicolumn{2}{|p{12.5cm}|}{能夠評估放射治療適應症,與放射腫瘤科協調治療計畫,理解放療技術,監測急慢性副作用。} \\
\hline
\textbf{核心能力} & 治療適應症 & 根治性、輔助性、緩和性放療決策 \\
\cline{2-3}
& 放療技術 & IMRT、SBRT、質子治療適應症理解 \\
\cline{2-3}
& 放化併用 & 同步vs序列放化療、藥物選擇 \\
\cline{2-3}
& 會診協調 & 與放射腫瘤科溝通、治療計畫整合 \\
\cline{2-3}
& 急性副作用 & 放射性皮膚炎、食道炎、腸炎處理 \\
\cline{2-3}
& 慢性副作用 & 纖維化、放射性肺炎、骨髓抑制 \\
\hline
\textbf{Level 1} & Novice & 能了解放療基本適應症 \\
\hline
\textbf{Level 2} & Advanced Beginner & 能評估基本放療需求,協調會診 \\
\hline
\textbf{Level 3} & Competent & 能獨立評估放療適應症,管理副作用 \\
\hline
\textbf{Level 4} & Proficient & 能協調複雜放化療,處理嚴重副作用 \\
\hline
\textbf{Level 5} & Expert & 成為放療協調專家,整合多模式治療 \\
\hline

\rowcolor{onco_teal!10}
\multicolumn{3}{|c|}{\textbf{EPA 7: 腫瘤急症診斷與處理}} \\
\hline
\textbf{描述} & \multicolumn{2}{|p{12.5cm}|}{能夠快速識別腫瘤急症,執行緊急評估與處置,包括SVCS、脊髓壓迫、高血鈣、TLS等,並能協調多科照護。} \\
\hline
\textbf{核心能力} & SVCS診斷處理 & 快速診斷、放療/支架決策 \\
\cline{2-3}
& 脊髓壓迫 & MRI評估、類固醇、放療/手術時機 \\
\cline{2-3}
& 高血鈣危象 & 快速識別、水化、雙磷酸鹽使用 \\
\cline{2-3}
& 腫瘤溶解症候群 & 預防、監測、Rasburicase使用 \\
\cline{2-3}
& 嗜中性白血球低下發燒 & 評估、經驗性抗生素、G-CSF使用 \\
\cline{2-3}
& 惡性腦壓 & 識別、類固醇、滲透壓治療 \\
\hline
\textbf{Level 1} & Novice & 能識別腫瘤急症症狀,知道基本處理 \\
\hline
\textbf{Level 2} & Advanced Beginner & 能執行初步評估,開始緊急處置 \\
\hline
\textbf{Level 3} & Competent & 能獨立處理腫瘤急症,協調多科照護 \\
\hline
\textbf{Level 4} & Proficient & 能處理複雜急症,快速決策與執行 \\
\hline
\textbf{Level 5} & Expert & 成為腫瘤急症專家,建立處理流程 \\
\hline

\rowcolor{onco_teal!10}
\multicolumn{3}{|c|}{\textbf{EPA 8: 化療副作用管理與支持性照護}} \\
\hline
\textbf{描述} & \multicolumn{2}{|p{12.5cm}|}{能夠預測、監測與處理化療副作用,執行支持性照護,包括止吐、G-CSF使用、黏膜炎處理、營養支持。} \\
\hline
\textbf{核心能力} & 噁心嘔吐 & 致吐風險評估、預防性止吐、救援治療 \\
\cline{2-3}
& 骨髓抑制 & 白血球、血小板低下監測與處理 \\
\cline{2-3}
& G-CSF使用 & 一級vs二級預防、使用時機 \\
\cline{2-3}
& 黏膜炎 & 口腔黏膜炎、腸道黏膜炎處理 \\
\cline{2-3}
& 周邊神經病變 & 評估、劑量調整、症狀緩解 \\
\cline{2-3}
& 營養支持 & 惡病質評估、營養介入、食慾促進劑 \\
\hline
\textbf{Level 1} & Novice & 能識別常見化療副作用 \\
\hline
\textbf{Level 2} & Advanced Beginner & 能處理常見副作用,執行支持治療 \\
\hline
\textbf{Level 3} & Competent & 能獨立管理化療副作用,預防性處置 \\
\hline
\textbf{Level 4} & Proficient & 能處理嚴重副作用,整合支持照護 \\
\hline
\textbf{Level 5} & Expert & 成為支持照護專家,改善副作用管理 \\
\hline

\rowcolor{onco_teal!10}
\multicolumn{3}{|c|}{\textbf{EPA 9: 癌症疼痛控制與緩和醫療}} \\
\hline
\textbf{描述} & \multicolumn{2}{|p{12.5cm}|}{能夠評估癌症疼痛,使用WHO三階梯止痛,處理難治性疼痛,執行症狀控制,並能提供全人緩和照護。} \\
\hline
\textbf{核心能力} & 疼痛評估 & VAS/NRS評分、疼痛特性、突發痛 \\
\cline{2-3}
& 鴉片類藥物 & 嗎啡、吩坦尼、羥考酮使用與轉換 \\
\cline{2-3}
& 輔助止痛 & 類固醇、抗痙攣藥、抗憂鬱劑 \\
\cline{2-3}
& 副作用管理 & 便秘、噁心、嗜睡處理 \\
\cline{2-3}
& 症狀控制 & 呼吸困難、譫妄、cachexia管理 \\
\cline{2-3}
& 心理社會支持 & 焦慮憂鬱、靈性照護、家屬支持 \\
\cline{2-3}
& 善終照護 & DNR討論、臨終症狀、哀傷輔導 \\
\hline
\textbf{Level 1} & Novice & 能評估基本疼痛,了解WHO階梯 \\
\hline
\textbf{Level 2} & Advanced Beginner & 能使用鴉片類藥物,處理常見副作用 \\
\hline
\textbf{Level 3} & Competent & 能獨立管理癌症疼痛,執行症狀控制 \\
\hline
\textbf{Level 4} & Proficient & 能處理難治性疼痛,提供全人照護 \\
\hline
\textbf{Level 5} & Expert & 成為緩和醫療專家,領導團隊照護 \\
\hline

\rowcolor{onco_teal!10}
\multicolumn{3}{|c|}{\textbf{EPA 10: 腫瘤病患跨科部整合照護}} \\
\hline
\textbf{描述} & \multicolumn{2}{|p{12.5cm}|}{能夠協調腫瘤病患的整體照護,包括MDT討論、跨科會診、臨床試驗評估、存活者照護,並能有效溝通與共同決策。} \\
\hline
\textbf{核心能力} & MDT參與 & 多專科團隊會議、治療計畫整合 \\
\cline{2-3}
& 會診協調 & 外科、放射、病理、影像科協調 \\
\cline{2-3}
& 臨床試驗 & 適應症評估、試驗選擇、納入標準 \\
\cline{2-3}
& 存活者照護 & 追蹤計畫、晚期效應、復發監測 \\
\cline{2-3}
& 病患溝通 & 壞消息告知、預後討論、共同決策 \\
\cline{2-3}
& 家屬支持 & 病情解釋、照護指導、資源連結 \\
\cline{2-3}
& 品質改善 & 照護品質監測、指標追蹤、改善方案 \\
\hline
\textbf{Level 1} & Novice & 能參與基本病患照護,了解團隊角色 \\
\hline
\textbf{Level 2} & Advanced Beginner & 能執行基本會診,參與MDT討論 \\
\hline
\textbf{Level 3} & Competent & 能獨立協調病患照護,進行有效溝通 \\
\hline
\textbf{Level 4} & Proficient & 能領導MDT討論,執行共同決策 \\
\hline
\textbf{Level 5} & Expert & 成為照護協調專家,改善照護品質 \\
\hline

\end{longtable}

\newpage
\section{評量方法與標準}

\subsection{CCC評量架構}
本腫瘤科EPAs採用Case-based Clinical Competency (CCC)評量方式:
\begin{itemize}[nosep]
    \item 真實臨床情境評量
    \item 多次觀察與回饋
    \item 漸進式委託決策
    \item 360度回饋機制
\end{itemize}

\begin{importantbox}
\textbf{特別提醒:} 腫瘤科EPAs的評量重點在於漸進式能力發展,而非單次完美表現。鼓勵從錯誤中學習,持續改進。每個EPA都需要在真實臨床情境中反復練習與評量。腫瘤學需要整合多學科知識,包括病理、影像、藥理、緩和醫療等。
\end{importantbox}

\subsection{評量工具對應表}
\begin{table}[h]
    \centering
    \caption{各EPA推薦評量工具}
    \begin{tabular}{p{4.5cm} p{4cm} p{5.5cm}}
        \toprule
        \textbf{EPA類型} & \textbf{主要評量工具} & \textbf{輔助評量方法} \\
        \midrule
        臨床推理類 & Mini-CEX & 病例討論、口試 \\
        (EPA 1, 3, 4, 5, 6, 7, 8) & (Mini-Clinical Evaluation Exercise) & 病歷撰寫評量、模擬情境 \\
        \midrule
        報告判讀類 & 判讀評量 & 線上測驗、定期考核 \\
        (EPA 2) & Report Interpretation Assessment & 病例討論、peer review \\
        \midrule
        緩和照護類 & OSCE/模擬病人 & 溝通技巧評量 \\
        (EPA 9) & Communication Skills Assessment & 360度回饋、錄影回顧 \\
        \midrule
        整合照護類 & Multi-source Feedback & MDT參與評量 \\
        (EPA 10) & (360-degree Assessment) & 病人滿意度調查 \\
        \bottomrule
    \end{tabular}
\end{table}

\subsection{評量頻率建議}
\begin{itemize}[nosep]
    \item \textbf{每月評量}:選擇3-4個EPA重點評量
    \item \textbf{季度檢討}:全面檢視所有10個EPA進展
    \item \textbf{年度認證}:EPA達標認證與進階規劃
    \item \textbf{即時回饋}:發現學習需求時立即介入
\end{itemize}

\section{實施建議與學習資源}

\subsection{月度晨會Mini-OSCE整合建議}
\begin{table}[h]
    \centering
    \caption{月度EPA評量主題安排}
    \begin{tabular}{p{2.5cm} p{4.5cm} p{7cm}}
        \toprule
        \textbf{月份} & \textbf{重點EPAs} & \textbf{評量重點} \\
        \midrule
        第1個月 & EPA 1, 2 & 基礎評估與病理判讀 \\
        第2個月 & EPA 3, 6 & 分期評估與放療協調 \\
        第3個月 & EPA 4, 8 & 化療執行與副作用管理 \\
        第4個月 & EPA 5 & 標靶與免疫治療 \\
        第5個月 & EPA 7, 9 & 急症處理與疼痛控制 \\
        第6個月 & EPA 10 & 整合照護與溝通能力 \\
        第7-12個月 & 綜合評量與進階訓練 & 複雜個案整合處理能力 \\
        \bottomrule
    \end{tabular}
\end{table}

\subsection{推薦學習資源}
學員可參考以下文獻深入學習:

\textbf{基礎教科書:}
\begin{itemize}
    \item DeVita, Hellman, and Rosenberg's Cancer: Principles \& Practice of Oncology. 11th edition.
    \item Abeloff's Clinical Oncology. 6th edition.
    \item Oxford Textbook of Oncology. 3rd edition.
\end{itemize}

\textbf{臨床指引:}
\begin{itemize}
    \item NCCN Clinical Practice Guidelines in Oncology
    \item ASCO Clinical Practice Guidelines
    \item ESMO Clinical Practice Guidelines
    \item WHO Guidelines for Cancer Pain Management
    \item MASCC Guidelines for Antiemetic Use
\end{itemize}

\textbf{線上資源:}
\begin{itemize}
    \item ASCO University
    \item ESMO e-Learning Platform
    \item UpToDate Oncology Section
    \item CancerNetwork
    \item Medscape Oncology
\end{itemize}

\subsection{自我評估工具}
\begin{itemize}[nosep]
    \item 定期記錄學習歷程與反思
    \item 建立個人病理報告判讀資料庫
    \item 參與MDT討論會
    \item 練習RECIST criteria判讀
    \item 建立個人學習檔案
    \item 參與腫瘤個案討論會
    \item 定期複習化療方案與劑量
    \item 練習壞消息告知技巧
\end{itemize}

\section{總結}

腫瘤科EPAs涵蓋了癌症診療的完整核心能力,從基礎的病史詢問與病理判讀,到複雜的化療管理、標靶免疫治療、腫瘤急症處理,以及全人化的緩和照護與跨科協調。每個EPA都有清楚的能力描述與里程碑發展路徑,協助住院醫師循序漸進地建立專業能力。

\begin{importantbox}
\textbf{關鍵訊息:} 腫瘤科EPAs的核心在於「可委託性」與「完整性」。這10個EPA代表腫瘤科醫師必須具備的基礎診療能力。腫瘤學需要整合病理、影像、藥理、緩和醫療等多學科知識。化療與標靶免疫治療需要精確評估適應症與監測副作用。緩和醫療與溝通技巧是腫瘤科醫師的核心能力,需要持續練習與反思。MDT討論與跨科協調能力對提供最佳癌症照護至關重要。
\end{importantbox}

% References section
\section{參考文獻}
\begin{thebibliography}{10}

\bibitem{devita2019}
DeVita, V. T., Lawrence, T. S., \& Rosenberg, S. A. (2019). \textit{DeVita, Hellman, and Rosenberg's Cancer: Principles \& Practice of Oncology}. 11th edition. Wolters Kluwer.

\bibitem{abeloff2020}
Niederhuber, J. E., et al. (2020). \textit{Abeloff's Clinical Oncology}. 6th edition. Elsevier.

\bibitem{olle2013}
ten Cate, O. (2013). Nuts and bolts of entrustable professional activities. \textit{Journal of Graduate Medical Education}, 5(1), 157-158.

\bibitem{who2018pain}
World Health Organization. (2018). \textit{WHO Guidelines for the Pharmacological and Radiotherapeutic Management of Cancer Pain in Adults and Adolescents}. WHO.

\bibitem{eisenhauer2009}
Eisenhauer, E. A., et al. (2009). New response evaluation criteria in solid tumours: revised RECIST guideline (version 1.1). \textit{European Journal of Cancer}, 45(2), 228-247.

\bibitem{roila2016}
Roila, F., et al. (2016). 2016 MASCC and ESMO guideline update for the prevention of chemotherapy- and radiotherapy-induced nausea and vomiting. \textit{Annals of Oncology}, 27(suppl 5), v119-v133.

\bibitem{brahmer2018}
Brahmer, J. R., et al. (2018). Management of immune-related adverse events in patients treated with immune checkpoint inhibitor therapy. \textit{Journal of Clinical Oncology}, 36(17), 1714-1768.

\end{thebibliography}

\end{document}